% ══════════════════════════════════════════════════════════════════
%  SECTION III --- BASELINE EVALUATION
% ══════════════════════════════════════════════════════════════════
\section{Baseline Evaluation}
\label{sec:baselines}
We use two kinds of baselines. The first is the five detectors (Table~\ref{tab:det}). The second is six explanation policies that the trigger is compared with in Section~\ref{sec:rq4}: explain nothing (B1); explain every flow with TreeSHAP or KernelSHAP (B2a/B2b); show one global explanation (B3); explain random flows at the trigger's rate (B4); explain every alert (B5); and use a learned fast explainer for every flow (B6).

\begin{table}[!htbp]
\centering
\caption{Detection results on the test set ($T_1$ = all features, $T_3$ = encrypted-traffic view). Alert threshold set on validation data for at most 1\% false positives; FPR in \%.}
\label{tab:det}
\scriptsize
\setlength{\tabcolsep}{2.5pt}
\begin{minipage}[t]{0.495\textwidth}
\centering
\begin{tabular}{llrrrrrr}
\toprule
\multicolumn{8}{c}{\textbf{CIC-IDS2017}} \\
\midrule
\textbf{Tier} & \textbf{Model} & \textbf{Prec.} & \textbf{Recall} & \textbf{F1} & \textbf{FPR} & \textbf{ROC} & \textbf{PR} \\
\midrule
T1 & LR & 0.994 & 0.815 & 0.896 & 0.5 & 0.991 & 0.987 \\
T1 & RF & 0.992 & 0.993 & 0.992 & 0.8 & 1.000 & 1.000 \\
T1 & XGB & 0.994 & 0.999 & 0.996 & 0.6 & 1.000 & 1.000 \\
T1 & MLP & 0.995 & 0.932 & 0.962 & 0.5 & 0.995 & 0.996 \\
\midrule
T3 & LR & 0.993 & 0.770 & 0.867 & 0.5 & 0.980 & 0.978 \\
T3 & EBM & 0.993 & 0.958 & 0.975 & 0.7 & 0.997 & 0.997 \\
T3 & RF & 0.992 & 0.984 & 0.988 & 0.8 & 0.998 & 0.998 \\
T3 & XGB & 0.993 & 0.969 & 0.981 & 0.7 & 0.998 & 0.998 \\
T3 & MLP & 0.994 & 0.928 & 0.960 & 0.5 & 0.993 & 0.994 \\
\bottomrule
\end{tabular}
\end{minipage}
\hfill
\begin{minipage}[t]{0.495\textwidth}
\centering
\begin{tabular}{llrrrrrr}
\toprule
\multicolumn{8}{c}{\textbf{HIKARI-2021}} \\
\midrule
\textbf{Tier} & \textbf{Model} & \textbf{Prec.} & \textbf{Recall} & \textbf{F1} & \textbf{FPR} & \textbf{ROC} & \textbf{PR} \\
\midrule
T1 & LR & 0.496 & 0.070 & 0.123 & 1.1 & 0.929 & 0.509 \\
T1 & RF & 0.497 & 0.155 & 0.236 & 2.4 & 0.930 & 0.513 \\
T1 & XGB & 0.492 & 0.123 & 0.197 & 1.9 & 0.930 & 0.504 \\
T1 & MLP & 0.554 & 0.097 & 0.165 & 1.2 & 0.937 & 0.554 \\
\midrule
T3 & LR & 0.495 & 0.067 & 0.118 & 1.0 & 0.929 & 0.509 \\
T3 & EBM & 0.502 & 0.159 & 0.242 & 2.4 & 0.932 & 0.519 \\
T3 & RF & 0.507 & 0.175 & 0.260 & 2.6 & 0.932 & 0.521 \\
T3 & XGB & 0.504 & 0.175 & 0.259 & 2.6 & 0.932 & 0.518 \\
T3 & MLP & 0.519 & 0.096 & 0.161 & 1.4 & 0.933 & 0.527 \\
\bottomrule
\end{tabular}
\end{minipage}
\end{table}


\begin{itemize}[leftmargin=*, itemsep=0pt]
  \item \textbf{CIC-IDS2017 is easy, even when encrypted.} The tree models and the MLP reach F1 0.96--0.996. Hiding ports and TCP headers ($T_1 \rightarrow T_3$) costs little: XGBoost 0.996 $\rightarrow$ 0.981 (precision 0.993, recall 0.969), Random Forest 0.992 $\rightarrow$ 0.988, MLP 0.962 $\rightarrow$ 0.960. Only LR loses more (0.896 $\rightarrow$ 0.867). The explainable EBM reaches 0.975~\cite{rudin2019}.
  \item \textbf{HIKARI-2021 is hard.} All models rank attacks well (ROC-AUC 0.93), but at 1\% false positives they catch only 7--18\% of attacks (F1 0.12--0.26), because HIKARI's normal traffic includes HTTPS logins to the same servers. Allowing 11--13\% false positives raises F1 to 0.69--0.73. HIKARI's explanation results are therefore read together with CIC-IDS2017 only.
  \item \textbf{Detection time per flow} (one core): 0.06--0.08\,ms for LR and the MLP, 0.36--0.55\,ms for the EBM and XGBoost, and 3.5--4.3\,ms for Random Forest. All overhead numbers in this report are relative to these times.
\end{itemize}
